\section{Benchmark}
\label{sec:benchmark}

\subsection{Datasets}
\paragraph{AMI (controlled, near-clean ground truth)} The AMI ear database~\cite{gonzalez2012ami} contains \NAmiAllImg{} images of 100 subjects, seven per subject, taken indoors at $492\times702$ pixels. The originals are stored at JPEG quality 90 and are downsized by a factor of 2.6 to 5 to a short side of 96, 144 and 192 pixels, which gives inputs of 24, 36 and 48 pixels at $\times4$, the size range of real small ear images. Downsizing by a non-integer factor removes the $8\times8$ block grid of the originals: the ratio of luminance differences across block boundaries to differences inside blocks is \BlockAmiRaw{} for the originals and \BlockAmiGtM{} for the 144-pixel ground truth (median over images; 1 means no grid). We therefore treat these downsized images as near-clean ground truth. The 36-pixel input is the main setting, and AMI is the dataset on which the fidelity results rest. Real small images have no ground truth; for them we report identification (Section~\ref{sec:recog}) and examples only. Subjects are split into five folds (70 for training, 10 for validation, 20 for testing). Two folds were held out during development and were trained and evaluated once, with a fixed configuration.

\paragraph{EarVN1.0 (in the wild)} EarVN1.0~\cite{hoang2019earvn} contains web-collected ear crops of \RoleAll{} subjects with widely varying size, pose and illumination. Each subject has one role, fixed before any experiment: \RoleTrain{} subjects supply additional training images, \RoleTest{} are used for testing, \RoleFit{} for measuring the degradation, \RoleClf{} for the realism classifier and \RoleViewer{} are reserved for evaluation on real small images. Ground truth is built from originals whose short side is at least twice the target size. Downsizing weakens but does not remove their compression (JPEG quality \GtEarvnQMin{} to \GtEarvnQMax{}, median \GtEarvnQMed{}), so PSNR on EarVN1.0 is measured against ground truth that is itself compressed. This yields \NEarvnImg{} test images of \NEarvnSubj{} subjects at 24-pixel inputs.

\paragraph{AWEx (in the wild, never used for training)} From the extended Annotated Web Ears collection~\cite{emersic2017ear,emersic2017uerc} we obtain, by the same rule, \NAwexSImg{} images of \NAwexSSubj{} subjects at 24-pixel inputs and \NAwexMImg{} images of \NAwexMSubj{} subjects at 36-pixel inputs. AWEx is distributed as PNG files whose compression history is not documented, and no AWEx image is used for training, fitting or model selection.

\subsection{Input degradations}
\label{sec:degradations}
All inputs are generated at $\times4$ with a fixed seed per image, so that all models receive identical inputs.
\emph{Clean}: bicubic downsampling, the NTIRE protocol.
\emph{JPEG $q$}: bicubic downsampling followed by JPEG compression at quality $q\in\{60,75,85,93\}$ (libjpeg through OpenCV, with 4:2:0 chroma subsampling).
\emph{Generic}: a single-stage pipeline in the style of~\cite{wang2021realesrgan,zhang2021bsrgan}, with Gaussian blur ($\sigma$ from 0.2 to 3.0 input pixels), additive Gaussian noise ($\sigma$ from 1 to 30 on the 8-bit scale) and JPEG compression (quality 30 to 90).
\emph{Measured}: the same structure with parameters measured from real small ear images (Section~\ref{sec:measured}).

\subsection{Models}
We evaluate every $\times4$ model with public weights that we could run without modification: \NPub{} fidelity-oriented models with distinct weights and four perceptual (GAN-trained) models, listed in Table~\ref{tab:published}. The fidelity-oriented set covers the NTIRE Efficient SR baseline SPAN~\cite{wan2024span} in three widths, entries of the 2023 to 2026 challenges~\cite{ren2024ntire,ren2025ntire,ren2026ntire}, RLFN~\cite{kong2022rlfn}, SAFMN++~\cite{sun2023safmn}, SMFANet~\cite{zheng2024smfanet}, and four larger references: MSRResNet~\cite{ledig2017srgan}, EDSR-baseline~\cite{lim2017edsr}, SwinIR-light~\cite{liang2021swinir} and the \ParamsRrdb{}M-parameter RRDB network~\cite{wang2018esrgan}. All were trained by their authors with bicubic downsampling. The perceptual set is ESRGAN~\cite{wang2018esrgan}, BSRGAN~\cite{zhang2021bsrgan} and two Real-ESRGAN models~\cite{wang2021realesrgan}.

\subsection{Metrics and statistics}
Fidelity is measured by PSNR and SSIM~\cite{wang2004ssim} on the luminance channel and perceptual quality by LPIPS~\cite{zhang2018lpips} and DISTS~\cite{ding2020dists}, with 32-bit inference. Comparisons are paired by image, and 95\% confidence intervals come from a bootstrap over subjects~\cite{efron1993bootstrap}. On AMI each fine-tuned model is evaluated on the test subjects of its own fold and the folds are pooled; on EarVN1.0 and AWEx the scores of the five fold models are averaged per image.
